\section{Background}
\label{sec:background}

The design of accountable, autonomous multi-agent systems requires a precise answer to a foundational question: when one agent delegates a task to another, through what structural mechanism is the chain of accountability preserved? This question becomes acute in systems that support \emph{recursive spawning} — the capacity for an agent to decompose its assigned task by provisioning child agents that operate with delegated authority. As agentic runtimes transition from single-agent to hierarchical multi-agent architectures \cite{wooldridge1995,stone2000}, the problem of managing the lifecycle of spawned agents, and of maintaining an unbroken accountability chain from every active agent back to an originating human decision, moves from an architectural footnote to a core safety requirement.

This section situates Recursive Spawning within five distinct research and engineering traditions: agent lifecycle management in multi-agent systems; accountability and provenance in distributed autonomous systems; fixed versus mutable delegation chains; hierarchical task decomposition in classical and modern AI; and the BX3 Framework's three-layer model as the architectural substrate for immutable parent chains. Together these traditions establish the conceptual and technical foundations upon which the Recursive Spawning Protocol (RSP), defined formally in Section~\ref{sec:rsp}, is constructed.

% -------------------------------------------------------
\subsection{Agent Lifecycle Management in Multi-Agent Systems}
% -------------------------------------------------------

Multi-agent systems (MAS) research has long recognized that the lifecycle of an agent — its creation, activation, operation, suspension, and termination — is a first-class design concern, not merely an implementation detail. The taxonomy introduced by Stone and Veloso \cite{stone2000} classifies multi-agent systems along axes of agent heterogeneity (homogeneous versus heterogeneous populations), agent knowledge (perfect versus imperfect information), agent access to the environment (fully versus partially accessible), and the nature of agent interactions (cooperative versus competitive). Within this taxonomy, recursive spawning is most relevant to cooperative, homogeneous systems operating in partially accessible environments — precisely the regime in which modern agentic runtimes such as AutoGPT \cite{autogpt2023}, LangChain \cite{langchain2023}, and Microsoft's AutoGen \cite{autogen2023} operate.

The management of agent lifecycles in such systems must address at minimum four distinct concerns: \emph{creation} (how an agent is instantiated and initialized with authority), \emph{delegation} (how authority is transmitted from parent to child), \emph{surveillance} (how the parent monitors the child's execution state), and \emph{termination} (how the child exits and returns results or an error state). Ferber and M\"uller \cite{ferber1996} formalized these as interactions between an agent's lifecycle state machine and the organizational context in which it operates, arguing that organizational roles — not just individual agent states — must be modeled to ensure system-wide coherence. This insight is directly relevant to Recursive Spawning: the spawn chain is not merely a tree of process instances; it is an organizational hierarchy of delegated authority, and each node's lifecycle must be governed with respect to both its local state and its position in that hierarchy.

The concept of agent delegation was formalized by Yim et al. \cite{yim1997} as a structured transfer of goals, plans, and execution authority from one agent to another, distinguished from simple message passing by the explicit surrender of control. Subsequent work by Padgham and Thiel \cite{padgham2001} demonstrated that incomplete lifecycle management — particularly the failure to distinguish between suspension and termination states — introduces subtle but systemic failures in agent systems, where orphaned or zombie agents continue to consume resources and emit effects long after their parent has exited. These findings establish that lifecycle management is not a single design decision but an ongoing architectural discipline: the system must maintain coherent authority boundaries at every state transition throughout the entire spawn chain.

More recent work on multi-agent orchestration has surfaced the specific problem of \emph{unbounded spawning}. The emergence of large language model-based agents capable of tool use and self-prompted task decomposition has made the uncontrolled proliferation of child agents a practical concern rather than a theoretical one. Zhuge et al. \cite{zhuge2024} document the ``agent chain problem'': LLM-based agents that recursively spawn children to handle sub-tasks, without a defined termination condition or depth ceiling, can produce chains of dozens of intermediate agents that collectively exceed human capacity for oversight and whose aggregate authority may differ substantially from what any single agent was authorized to exercise. This problem is not addressed by classical lifecycle management alone, because the agent chain in such systems is implicit — encoded in context windows and prompt structures rather than in first-class authorization artifacts — and therefore not amenable to systematic audit.

% -------------------------------------------------------
\subsection{Accountability and Provenance in Distributed Autonomous Systems}
% -------------------------------------------------------

The accountability problem in distributed autonomous systems has been studied under several related formalisms. The Human Delegation Provenance (HDP) Protocol represents one of the most rigorous approaches: a lightweight, cryptographic protocol that binds human authorization to a session and then records every subsequent agent delegation as a signed, append-only hop in a provenance chain \cite{hdp2026,hdp2025ietf}. Each hop is cryptographically signed using Ed25519, and the complete chain can be verified offline using only the issuer's public key and the session identifier — no central registry or trust anchor is required. The HDP token structure encodes header, principal, scope, chain of custody, and signatures, ensuring that the scope of delegated authority is itself an auditable artifact \cite{hdp2026}. This combination of scope-enriched delegation tokens and offline-verifiable chain-of-custody directly addresses the provenance requirement for Recursive Spawning: every spawned child must carry, as a first-class artifact, a verifiable chain back to the human who authorized its parent's authority.

Chadwick and Kaiser \cite{chadwick2023} extend this line of work by examining the specific failure modes that arise when delegation chains grow long and the original human principal is no longer identifiable by downstream agents. Their ``authorization drift'' phenomenon — in which agents progressively replace human intent with locally optimized goals — is structurally analogous to what we term \emph{autonomous drift} in the Recursive Spawning context (Section~\ref{sec:drift}). Both phenomena share a common root: the absence of a structurally enforced human anchor at the terminus of the delegation chain. Chadwick and Kaiser's mitigation strategy relies on periodic re-authorization by the originating human; our mitigation relies on the immutable parent pointer and the upstream accountability guarantee provided by the BX3 Framework's Fact Layer, which together ensure that every active agent has a traversable, unbreakable chain to a human principal at all times — not just at authorization time.

Xu et al. \cite{xu2023} formalize accountability in distributed systems using provenance graphs, demonstrating that accountability guarantees are only as strong as the immutability guarantees of the provenance record itself. If the provenance graph is mutable — if edges can be rewritten after the fact — then any accountability guarantee derived from it is conditional on the absence of a retroactive edit. This motivates the architectural requirement that the parent pointer chain must be immutable from the moment of spawn: the warrant that activates a child agent must encode an immutable delegation chain, because a mutable chain provides no more than a probabilistic accountability signal.

% -------------------------------------------------------
\subsection{Fixed Versus Mutable Delegation Chains}
% -------------------------------------------------------

Delegation chains in multi-agent systems take two structurally different forms with profoundly different accountability properties. In a \emph{fixed delegation chain}, the parent-child relationship is established at spawn time and is immutable thereafter: the child's parent pointer cannot be changed, the child's authority cannot be re-delegated except through the defined spawn protocol, and the chain terminates at a predetermined anchor (typically a human principal). In a \emph{mutable delegation chain}, the parent-child relationship is treated as runtime state that can be updated, redirected, or reassigned as conditions change.

Mutable delegation chains offer greater flexibility: a child agent can be reparented to a different parent if the organizational structure changes, if a parent fails, or if load balancing requires reassignment. This flexibility, however, comes at a significant accountability cost. If the parent-child relationship is mutable at runtime, then any single action by a child agent may have been authorized by one parent at one moment and a different parent at another — and the effective authorization for any given action may not be the authorization that was recorded at spawn time. The chain of provenance becomes conditional on the state of a mutable mapping rather than fixed by an immutable record. This introduces what we term the \emph{delegatee ambiguity problem}: after an event, it may be impossible to determine which principal effectively authorized the action, because the delegation chain has been rewritten since the action occurred.

Fixed delegation chains sacrifice this flexibility for a stronger accountability guarantee. Each child's authority is traceable to a single, immutable parent pointer, and that pointer chains unbroken back to a human anchor. There is no ambiguity about effective authorization: the warrant that activated the child encodes the complete delegation chain as it existed at the moment of activation, and that chain cannot be altered retroactively. The cost is that reparenting requires a formal re-spawn: the child must be deactivated and a new child instantiated under the new parent, with a new warrant and a new immutable chain. This cost is architectural, not practical — it is the price of making accountability structurally derivable rather than forensically reconstructed.

The contrast between fixed and mutable chains maps directly onto the distinction between static and dynamic organizational structures in MAS theory. Fox \cite{fox1981} and later Spector and Gopal \cite{spector1994} demonstrate that static organizational structures provide superior traceability in cooperative multi-agent systems but suffer from brittleness under environmental change. The BX3 Framework resolves this tension not by choosing one structure over the other but by layering them: the spawn chain is fixed and immutable (providing traceability and accountability), while the Fact Layer's forensic ledger records the full operational history of each node (providing adaptability and auditability) without requiring the chain itself to be mutable.

% -------------------------------------------------------
\subsection{Hierarchical Task Decomposition in Classical and Modern AI}
% -------------------------------------------------------

Hierarchical task decomposition is a foundational technique in AI planning, stretching from classical hierarchical task network (HTN) planning \cite{erol1994,ghallab2004} through modern LLM-based agent architectures \cite{wei2022,shinn2023}. In classical HTN planning, a complex task is decomposed recursively into simpler subtasks, each of which is assigned to a distinct planning node. The decomposition is planned in advance by the planner before execution begins; the plan is a tree of tasks, and execution proceeds by selecting methods for each task and expanding the tree until all leaves are primitive actions \cite{ghallab2004}. Classical HTN planning assumes a single, central planner with perfect knowledge of the task and its environment — an assumption that does not hold in distributed autonomous systems where multiple agents operate concurrently with partial information.

Modern LLM-based agents adopt a variant of hierarchical decomposition in which the agent itself decides, at runtime, when and how to decompose a task based on encountered subtask complexity or failed actions. This dynamic decomposition is enabled by the model's capacity for in-context reasoning: an LLM agent can observe that a current action has failed, infer that the failure stems from task complexity beyond its provisioned scope, and autonomously spawn a child agent to handle a specific sub-task \cite{press2022,yao2022}. Unlike classical HTN planning, the decomposition is not pre-planned; it is reactive and context-sensitive. The spawn tree grows dynamically based on the agent's real-time observations, allowing the agent to adapt its decomposition as task conditions change.

What connects classical HTN and modern LLM-based decomposition is the spawn tree: a hierarchical structure in which each node represents a subtask and edges represent the decomposition relationship. In both traditions, the tree is typically mutable and implicit — decomposition decisions are not recorded as first-class authorization artifacts, and child nodes are not distinguished from parent nodes in terms of their accountability status. A classical HTN planner's subtasks are subroutine calls; they do not have independent accountability chains. An LLM agent's spawned sub-agents are independent processes, but the chain of authorization is implicit in the agent's context window rather than explicit in an immutable warrant structure.

Recursive Spawning adopts the hierarchical decomposition paradigm from both traditions but adds the critical accountability layer that is absent from both: each node in the decomposition tree is a first-class accountability entity with an immutable parent pointer, a defined termination condition, and a cryptographically verifiable chain back to a human anchor. Decomposition is not merely a planning technique; it is an authorization event that creates a new accountable entity. The decomposition tree is therefore not just a planning artifact but an organizational chart of delegated authority.

% -------------------------------------------------------
\subsection{The BX3 Framework as Substrate for Immutable Parent Chains}
% -------------------------------------------------------

The BX3 Framework provides the architectural substrate within which Recursive Spawning operates. As defined in \cite{bx3framework2026}, the BX3 Framework organizes any autonomous node into three immutable functional layers — the \textit{Purpose Layer}, the \textit{Bounds Engine}, and the \textit{Fact Layer} — each defined by the properties it must maintain rather than by the type of actor occupying it. The Purpose Layer carries intent, judgment, and final human accountability; the Bounds Engine carries constrained execution within a defined operating range; the Fact Layer carries deterministic enforcement and forensic auditability.

The upstream accountability guarantee of the BX3 Framework states that when any node encounters a state it cannot resolve within its provisioned bounds, accountability escalates recursively upward through the system hierarchy, bypassing all machine actors, until it reaches a human anchor \cite{bx3framework2026}. This guarantee is not a policy choice; it is a structural property of the three-layer architecture. Specifically, the Fact Layer creates the immutable parent pointer at spawn time, issues the spawn warrant, and maintains the forensic ledger that records every spawn event. The Fact Layer's determinism is what makes the parent pointer immutable: once written, the pointer cannot be updated because the Fact Layer does not permit modifications to its authorization records.

The three-layer model provides the structural basis for the upstream accountability guarantee in the spawn context specifically. When a child agent is activated via a spawn warrant, its authority is bounded by the operating range defined by the parent node's Bounds Engine. If the child encounters a condition outside its provisioned scope, it triggers the Bailout Protocol \cite{bailoutprotocol2026} — the upward exception propagation mechanism of the BX3 Framework — which bypasses all intermediate machine actors and routes the exception to the nearest human accountability anchor. The spawn chain is therefore not only an authorization structure but an escalation structure: the same chain that carries delegated authority from the human downward carries exceptions upward when the Bounds Engine cannot resolve them.

This coupling of authorization and escalation is the key architectural insight of Recursive Spawning. In conventional systems, authorization and escalation are separate concerns: a spawning decision is made by the parent agent, and an escalation decision is made independently when an exception occurs. In the BX3 Framework, both decisions are governed by the same immutable chain structure. There is no separate escalation infrastructure; the spawn chain is the escalation chain, by design. The result is an accountability architecture in which the question — \emph{who is responsible for this agent's actions?} — is always answerable by traversing the immutable parent chain back to a human anchor, and in which every traversal of that chain is a structurally enforced property of the runtime rather than a forensic reconstruction performed after the fact.